\begin{lemma}[From weak profile convergence to empirical hypotheses]
\label{lem:random-profile-bridge}
Fix $T>0$, put $L=\log N$, $K=\lfloor TL\rfloor$, and
$S_N=K/L$, and write $p^{(j)}=X_N^{(j)}/N$ for the random pile
proportions. Suppose that the weak convergence in probability
in~\eqref{eq:randomized-profile-like} holds at horizon $T$, namely
\[
 M_N^T=\frac1L\sum_{j=1}^K\delta_{(p^{(j)},j/L)}
 \xrightarrow{\mathrm{weakly,\ probability}}
 M^T(dp,dt)=\mu_t(dp)\,dt
 \quad\text{on }D_k\times[0,T].
\]
Then the random error
\begin{align*}
 e_N=\max\bigg\{&
 \sup_{0\le r\le s\le S_N}|A^{(N)}(r,s)-A(r,s)|,\\
 &\sup_{\substack{0\le r\le s\le S_N\\2\le\theta\le4}}
       |B_\theta^{(N)}(r,s)-B_\theta(r,s)|\bigg\}
\end{align*}
converges to zero in probability. In particular,
\begin{equation}\label{eq:random-profile-margin-convergence}
 \beta^{(N)}(S_N)\xrightarrow{\mathbb P}\Lambda_0(T),
 \qquad
 \Lambda_1^{(N)}\xrightarrow{\mathbb P}\Lambda_1(T).
\end{equation}
If, in addition, the nondegeneracy condition
in~\eqref{eq:profile-nondegeneracy} holds on $[0,T]$, so that
\[
 a(t):=\int_{D_k}a_p\,d\mu_t(p)\ge a_T>0
 \quad\text{for almost every }t\in[0,T],
\]
there are fixed $g\in(0,1/k)$ and $\gamma>0$, depending only on
$k,a_T$, such that, for every fixed $\rho>0$,
\begin{equation}\label{eq:random-profile-local-good}
 \Pp\{\text{the array is locally }(g,\gamma,\rho)
                  \text{-good on }[0,T]\}\longrightarrow1.
\end{equation}
The same $g,\gamma$ work for every $\rho$; the convergence is asserted
separately for each fixed $\rho$.
\end{lemma}

\begin{proof}
For a nonnegative continuous function $f$ on $D_k$, define
\[
 F_N^f(t)=\frac1L\sum_{\substack{1\le j\le K\\j/L\le t}}f(p^{(j)}),
 \qquad
 F^f(t)=\int_0^t\int_{D_k}f(p)\,d\mu_v(p)\,dv
 \quad(0\le t\le T).
\]
The limiting measure $M^T$ gives zero mass to every fixed time slice.
Consequently weak convergence in probability gives
$F_N^f(t)\to F^f(t)$ in probability at every fixed $t$.
For example, this follows by bounding the indicator of $[0,t]$
above and below by continuous functions of time; the limiting error
is bounded by $\|f\|_\infty$ times the lengths of their transition
intervals.

Both cumulative functions are nondecreasing, and $F^f$ is
$\|f\|_\infty$-Lipschitz. For a deterministic mesh
$0=t_0<\cdots<t_m=T$ with gaps at most $\delta$, monotonicity gives
\[
 \sup_{0\le t\le T}|F_N^f(t)-F^f(t)|
 \le\max_{0\le i\le m}|F_N^f(t_i)-F^f(t_i)|
                         +\|f\|_\infty\delta.
\]
First use convergence on this finite mesh and then let $\delta$
decrease to zero. This proves uniform convergence in probability
for each fixed $f$.

Apply this argument to $a$ and to $\psi_\theta$, where
$\psi_\theta(p)=\psi_p(\theta)$. Joint continuity on
$D_k\times[2,4]$ makes the family $\{\psi_\theta:2\le\theta\le4\}$
compact in the uniform norm. A finite net in this family suffices:
if $\|f-f'\|_\infty\le\varepsilon$, then, since
$M_N^T(D_k\times[0,T])=S_N\le T$,
\[
 \sup_{0\le t\le T}|F_N^f(t)-F_N^{f'}(t)|\le T\varepsilon,
 \qquad
 \sup_{0\le t\le T}|F^f(t)-F^{f'}(t)|\le T\varepsilon.
\]
It follows that
\[
 D_N:=\max\left\{
 \sup_{0\le t\le T}|F_N^a(t)-F^a(t)|,
 \sup_{\substack{0\le t\le T\\2\le\theta\le4}}
       |F_N^{\psi_\theta}(t)-F^{\psi_\theta}(t)|
                    \right\}\xrightarrow{\mathbb P}0.
\]

For $0\le t\le S_N$, the interpolated empirical integral of $f$
over $[0,t]$ differs from $F_N^f(t)$ by at most
$\|f\|_\infty/L$: only the column containing $t$ can contribute
partially. Taking differences at two endpoints and using
$\|a\|_\infty,\|\psi_\theta\|_\infty\le4\log k$ gives
\[
 e_N\le2D_N+\frac{8\log k}{L}
          \xrightarrow{\mathbb P}0.
\]
The deterministic comparison~\eqref{eq:margin-transfer}, including
the rounding error $0\le T-S_N<1/L$, now gives
\[
 |\beta^{(N)}(S_N)-\Lambda_0(T)|
       \le e_N+\frac{\log k}{L},\qquad
 |\Lambda_1^{(N)}-\Lambda_1(T)|
       \le\frac{13}{2}e_N+\frac{14\log k}{L}.
\]
The measurability of the profile ensures the scalar integrability
needed for that comparison, since all the scalar costs are bounded.
This proves~\eqref{eq:random-profile-margin-convergence}.

To prove local goodness, choose $g\in(0,1/k)$ sufficiently small
that
\[
 -\log(1-(k-1)g)\le a_T/4,
 \qquad \gamma=\frac{a_T}{8\log k}.
\]
If a column is not $g$-good, every coordinate except a maximizing
one is less than $g$. Thus its maximal coordinate is greater than
$1-(k-1)g$, and its cost $a_p$ is at most $a_T/4$.
All columns, good or otherwise, have cost at most $\log k$.

Fix $0<\rho\le T$. For any real $a$ with
$[a,a+\rho]\subseteq[0,T]$, let $G_N(a)$ count the $g$-good
columns whose times belong to this closed interval. The uniform
cumulative estimate and nondegeneracy give
\[
 \frac1L\sum_{\substack{1\le j\le K\\a\le j/L\le a+\rho}}
         a_{p^{(j)}}
 \ge F_N^a(a+\rho)-F_N^a(a)
 \ge a_T\rho-2D_N.
\]
There are at most $\rho L+1$ column times in the closed interval.
The bounds on good and nongood costs therefore imply
\[
 \frac{G_N(a)}L
 \ge\frac{3a_T\rho/4-2D_N-a_T/(4L)}{\log k}.
\]
This holds simultaneously for every real $a$ in the indicated range,
including intervals meeting the final gap $(S_N,T]$. On the event
$D_N\le a_T\rho/8$, and for $L^{-1}\le\rho$, its right side is
at least $a_T\rho/(4\log k)\ge\gamma\rho$.
The event has probability tending to one, proving
\eqref{eq:random-profile-local-good}. For $\rho>T$ the local
condition is vacuous.
\end{proof}
